\subsection{R 的势}

我们有 \(\mathbf{R} = \bigcup_{n\in \mathbf{Z}}[n,n + 1[,\) 且所有区间 \([n,n + 1[\) 都与 \([0,1[\) 等势。由于 \([0,1[\) 是无限集，由此推出 \(\mathbf{R}\) 与区间 \([0,1[\) 等势。通过考虑区间 \([0,1[\) 中数的二进展开，我们将证明这个区间与由每一项都等于 0 或 1 的序列 \((u_n)\) 组成的集合 S 等势。

首先，它与子集 \(S'\) （\(S\) 的，由序列 \((u_n)\) （满足 \(u_n = 0\) ，对无穷多个 \(n\) 的值）组成）等势。另一方面，余集 \(S''\) （\(S'\) 在 \(S\) 中的）与等于分母为 \(2^n\) 的分数的有理数的所有非正常展开组成的集合等势；这些数构成 \(Q\) 的一个子集，因此它们组成的集合是可数的，从而 \(S''\) 也是可数的。由于 \(S'\) 是无限的，它与 \(S\) 等势，由此即得结果。

现在 S 与 \(\mathfrak{P}(\mathbf{N})\) 等势；因为我们可以定义 \(\mathfrak{P}(\mathbf{N})\) 到 S 上的一个双射，把 N 的每个子集 A 映为序列 \((u_n)\) ，其中 \(u_n = 0\) 若 \(n \in A\) ，而 \(u_n = 1\) 若 \(n \in \mathbb{C}A\) 。我们这样就证明了：

定理 1（Cantor）. 实数组成的集合与一个可数无限集的一切子集组成的集合等势。

推论. 实数组成的集合的势严格大于可数集的势。

与 R 等势的集合称为具有连续统的势。由 § 4 第 1 节命题 1，每个含有多于一个点的区间都具有连续统的势。同样，R 的可数子集的余集具有连续统的势；特别地，所有无理数组成的集合具有连续统的势。

